\documentclass[12pt]{article}

\usepackage[utf8]{inputenc}
\usepackage{amsmath, amssymb, amsthm}
\usepackage{geometry}
\geometry{a4paper, margin=2.5cm}

\title{Decompoziția unui Modul Finit Generat peste un PID}
\author{}
\date{}

\theoremstyle{definition}
\newtheorem{definition}{Definiție}

\theoremstyle{plain}
\newtheorem{theorem}{Teoremă}

\begin{document}

\maketitle

\begin{definition}
Fie $R$ un domeniu principal (PID) și $M$ un $R$-modul. Partea torsională a lui $M$ este


\[
t(M) = \{\, x \in M : \exists\, 0 \ne r \in R \text{ cu } rx = 0 \,\}.
\]


\end{definition}

\begin{theorem}[Decompoziția torsională]
Fie $R$ un PID și $M$ un $R$-modul finit generat. Atunci există un submodul liber
$F \subseteq M$ astfel încât


\[
M = t(M) \oplus F.
\]


În special, $F$ este un modul liber de rang egal cu rangul lui $M$.
\end{theorem}

\begin{theorem}[Unicitatea părții libere]
În ipotezele teoremei precedente, rangul submodulului liber $F$ este unic, iar orice două
submodule libere $F_1, F_2 \subseteq M$ pentru care


\[
M = t(M) \oplus F_i
\]


sunt izomorfe ca $R$-module libere.
\end{theorem}

\end{document}
